\section{总结与延伸}

\subsection{核心内容总表}
\begin{table}[H]
\centering
\begin{tabular}{p{2.8cm}p{4.5cm}p{4.2cm}}
\toprule
\textbf{主题} & \textbf{课程结论} & \textbf{对 LLM Agent 启发} \\
\midrule
Agent 定义 & 连续感知-决策-行动闭环 & 评测要覆盖多步行为结果 \\
智能体架构 & Perception-Cognition-Action 分层 & 系统可调试性优先于端到端神化 \\
多体互动 & 协作与对抗是核心难题 & 多 Agent 场景需建模队友分布变化 \\
Human feedback & TAMER/RLHF 可加速早期学习 & 人类偏好与环境反馈需联合优化 \\
Embodied RL & Slack 通过结构化动作空间降本提效 & 高层 action primitive 可迁移到软件 Agent \\
竞争型 RL & GT Sophy 依赖系统工程和数据策划 & replay/curriculum 设计是策略质量关键 \\
\bottomrule
\end{tabular}
\caption{Autonomous Agents 课程要点浓缩}
\end{table}

\subsection{进一步思考}
\begin{enumerate}
    \item 当 Agent 要在真实业务里长期运行时，哪些失败模式必须在训练阶段提前暴露？
    \item 如何把 ``ad hoc teamwork'' 形式化到今天的 multi-agent coding pipeline？
    \item 在 LLM Agent 里，哪些能力应留在系统外层，哪些应内化到策略模型？
\end{enumerate}

\subsection{拓展阅读}
\begin{itemize}
    \item Peter Stone 个人主页与公开讲座（含 Autonomous Agents、GT Sophy、robot learning 系列）
    \item TAMER / Deep TAMER / human-in-the-loop reinforcement learning 相关论文
    \item Ad Hoc Teamwork challenge paper（AAAI 2010）与后续多智能体协作研究
    \item Slack（Sim-to-real latent action learning）与 factorized RL 相关论文
    \item Outracing champion Gran Turismo drivers with deep RL（Nature 2022）
    \item Phoebe dataset（human-centric fairness benchmark，Nature 2025）
\end{itemize}

\end{document}
